\chapter{Intertemporal Choice}\label{ch:b2:intertemporal-choice}

Income comes in a lump in mid-life: little at the start of a career, the most in the forties and
fifties, much less in retirement. Spending need not follow. A household can borrow when young, save
in its prime and draw on its savings when old, and the theory of this chapter says that a patient
household facing a fair interest rate will spread its consumption much more evenly than its income.
The chart puts the two side by side for the United States, one age group at a time.

\begin{openingchart}{Income after taxes and spending of American households by age of the reference
  person, 2023, in thousands of dollars a year: one point per age group, placed at the group's mean
  age. Spending is the survey's average annual expenditures less personal insurance and pensions
  (pension contributions are saving and Social Security contributions a tax: neither is consumption).}{U.S. Bureau of Labor
  Statistics, Consumer Expenditure Surveys, Table 1300 (age of reference person), 2023}
\begin{tikzpicture}
\begin{axis}[omrecord, xmin=15, xmax=88, ymin=0, ymax=130, xtick={20,30,...,80},
    xlabel={Mean age of the reference person (years)}, ylabel={Thousands of dollars a year},
    legend pos=north east, height=0.55\linewidth, width=0.8\linewidth]
  \addplot[omDef, very thick, mark=*, mark size=1.8pt] table[x=age, y=income_after]
    {figdata/bachelor-2/03-intertemporal-choice/out/ce_age.csv};
  \addplot[omThm, very thick, mark=square*, mark size=1.8pt] table[x=age, y=spending]
    {figdata/bachelor-2/03-intertemporal-choice/out/ce_age.csv};
  \legend{Income after taxes, Spending}
\end{axis}
\end{tikzpicture}
\end{openingchart}

Spending does rise less than income: from the youngest group to the 45--54 group, income after taxes
% figdata: ce_summary (peak_income_ratio 2.26, peak_spending_ratio 1.85)
is multiplied by 2.26 and spending by 1.85. In the oldest group, spending is 102 per cent of income,
% figdata: ce_age (ratio 1.023 for 75 and older)
against 71 per cent at 35--54. But spending still rises and falls with income, far more than a flat
% figdata: ce_age (ratio 0.713)
plan would allow. This chapter builds the model that predicts the flat plan, finds what bends it
(impatience, the interest rate, limits on borrowing), and tests a second model, in which people are
impatient in a way that changes as the date of a choice approaches.

\begin{recall}
From the Year 1 volume: the two-period budget constraint, saving as the difference between income and
consumption, the real interest rate as the price of consumption today in terms of consumption
tomorrow. From the school volumes: present value, $\sum_t m_t/(1+r)^t$. From
\cref{ch:b2:expected-utility}: CRRA utility, $u(c) = c^{1-\gamma}/(1-\gamma)$, or $\ln c$ when $\gamma
= 1$.
\end{recall}

\begin{notation}[Symbols of this chapter]\label{not:b2:intertemporal-choice:symbols}
Periods are $t = 0, 1, \dots, T - 1$; in period $t$ a household earns income $m_t$, consumes $c_t$
and ends the period with assets $b_t$ (negative when it owes). The real interest rate $r$ is the same
in every period. $\beta$ is the \omterm{def:b2:intertemporal-choice:discount-factor}{discount factor}, $\rho$ the \omterm{def:b2:intertemporal-choice:time-preference}{rate of time preference}, $\gamma$ the CRRA
coefficient, $\mathcal{M}$ the present value of lifetime income, $\bar B$ a borrowing limit, and
$\varphi$ the present-bias factor of \cref{sec:b2:intertemporal-choice:consistency}.
\end{notation}

\section{Discounting and time-separable utility}\label{sec:b2:intertemporal-choice:discounting}

A household chooses a whole plan of consumption $(c_0, c_1, \dots, c_{T-1})$. The simplest way to
rank plans adds up the utility of each period, with less weight on later ones.

\begin{definition}[Discount factor]\label{def:b2:intertemporal-choice:discount-factor}
A household with \emph{exponential discounting} ranks consumption plans by
\[
U_0 = \sum_{t=0}^{T-1} \beta^{t}\, u(c_t), \qquad 0 < \beta \le 1,
\]
where $u$ is increasing and concave. The number $\beta$ is its \emph{discount
factor}\index{discount factor}: one unit of utility a period later is worth $\beta$ units now.
\end{definition}

\begin{definition}[Rate of time preference]\label{def:b2:intertemporal-choice:time-preference}
The \emph{rate of time preference}\index{rate of time preference} $\rho$ is defined by $\beta = 1/(1 +
\rho)$. It is the rate at which the household discounts utility, as $r$ is the rate at which the market
discounts money.
\end{definition}

The model makes two strong assumptions. Utility is \emph{time-separable}: what a household gets from
consumption this year does not depend on what it consumed last year (no habits). And the weight on a
period depends only on its distance from today, through the same factor $\beta$ for each extra period.
\Cref{sec:b2:intertemporal-choice:consistency} shows what the second assumption buys.

The household can borrow and lend at the rate $r$. Starting with no assets, its assets evolve as
\[
b_t = (1 + r)\, b_{t-1} + m_t - c_t, \qquad b_{-1} = 0,
\]
and it may not die in debt: $b_{T-1} \ge 0$.

\begin{proposition}[Intertemporal budget constraint]\label{prop:b2:intertemporal-choice:budget}
A plan satisfies the period constraints with $b_{T-1} = 0$ if and only if
\[
\sum_{t=0}^{T-1} \frac{c_t}{(1 + r)^{t}} = \sum_{t=0}^{T-1} \frac{m_t}{(1 + r)^{t}} = \mathcal{M}.
\]
The present value of consumption equals the present value of income. With $b_{T-1} \ge 0$ the left
side is at most $\mathcal{M}$, and a household whose utility rises with consumption leaves nothing.
\end{proposition}

\begin{proof}
Divide the period constraint by $(1 + r)^{t}$:
$b_t/(1 + r)^{t} - b_{t-1}/(1 + r)^{t-1} = (m_t - c_t)/(1 + r)^{t}$. Summing over $t = 0, \dots, T - 1$,
the left side telescopes to $b_{T-1}/(1 + r)^{T-1}$, so $b_{T-1} \ge 0$ is $\sum_t (m_t - c_t)/(1 +
r)^{t} \ge 0$, with equality exactly when $b_{T-1} = 0$. If $b_{T-1} > 0$, raising $c_{T-1}$ by
$b_{T-1}$ keeps every constraint and raises utility.
\end{proof}

So the timing of income matters only through $\mathcal{M}$, the household's lifetime resources: two
income profiles with the same present value allow the same consumption plans, as long as the household
can borrow freely. \Cref{sec:b2:intertemporal-choice:borrowing} drops that condition.

\section{The Euler equation}\label{sec:b2:intertemporal-choice:euler}

\begin{definition}[Euler equation]\label{def:b2:intertemporal-choice:euler-equation}
The \emph{Euler equation}\index{Euler equation} of a consumption plan is the condition, for each pair
of adjacent periods,
\[
u'(c_t) = \beta\,(1 + r)\, u'(c_{t+1}).
\]
Giving up one unit of consumption at $t$ costs $u'(c_t)$; saved, it returns $1 + r$ units at $t + 1$,
worth $\beta(1 + r)u'(c_{t+1})$ in today's utility. At an optimum the two are equal.
\end{definition}

\begin{proposition}[The optimal plan]\label{prop:b2:intertemporal-choice:euler}
Let $u$ be strictly increasing, strictly concave and differentiable, with $u'(c) \to \infty$ as $c \to
0$. A plan maximises $U_0$ subject to the intertemporal budget constraint if and only if it meets the
constraint with equality and satisfies the \omterm{def:b2:intertemporal-choice:euler-equation}{Euler equation} for $t = 0, \dots, T - 2$.
\end{proposition}

\begin{proof}
The Lagrangian is $\sum_t \beta^{t} u(c_t) + \lambda\bigl(\mathcal{M} - \sum_t c_t/(1 + r)^{t}\bigr)$.
Since $u'$ is unbounded near zero, every $c_t$ is positive at the optimum, and the first-order
conditions are $\beta^{t} u'(c_t) = \lambda(1 + r)^{-t}$ for every $t$. Dividing the condition for $t +
1$ by the one for $t$ gives $\beta\,u'(c_{t+1})/u'(c_t) = 1/(1 + r)$, the \omterm{def:b2:intertemporal-choice:euler-equation}{Euler equation}. Conversely,
the objective is strictly concave and the constraint linear, so a plan that meets the constraint and
the first-order conditions (with $\lambda = u'(c_0)$) is the unique maximum.
\end{proof}

The \omterm{def:b2:intertemporal-choice:euler-equation}{Euler equation} fixes the \emph{slope} of the plan, the growth of consumption from one period to the
next; the budget fixes its \emph{level}. With CRRA utility the slope has a closed form.

\begin{corollary}[Consumption growth under CRRA]\label{cor:b2:intertemporal-choice:crra-growth}
If $u(c) = c^{1-\gamma}/(1 - \gamma)$ (or $\ln c$), then
\[
\frac{c_{t+1}}{c_t} = \bigl(\beta(1 + r)\bigr)^{1/\gamma}, \qquad \ln\frac{c_{t+1}}{c_t} =
\frac{\ln(1 + r) - \ln(1 + \rho)}{\gamma} \approx \frac{r - \rho}{\gamma}.
\]
Consumption grows when the market rate exceeds the \omterm{def:b2:intertemporal-choice:time-preference}{rate of time preference}, falls when it is below,
and grows faster the smaller is $\gamma$.
\end{corollary}

\begin{proof}
$u'(c) = c^{-\gamma}$, so the \omterm{def:b2:intertemporal-choice:euler-equation}{Euler equation} reads $c_t^{-\gamma} = \beta(1 + r)c_{t+1}^{-\gamma}$; take
the power $1/\gamma$ and logarithms, and use $\ln(1 + x) \approx x$ for small $x$.
\end{proof}

\begin{definition}[Elasticity of intertemporal substitution]\label{def:b2:intertemporal-choice:eis}
The \emph{elasticity of intertemporal substitution}\index{elasticity of intertemporal substitution}
(EIS) is the percentage change in the ratio $c_{t+1}/c_t$ caused by a one per cent rise in the gross
interest rate $1 + r$, along the \omterm{def:b2:intertemporal-choice:euler-equation}{Euler equation}. By \cref{cor:b2:intertemporal-choice:crra-growth} it
equals $1/\gamma$ under CRRA utility.
\end{definition}

So CRRA ties together two quite different things: the aversion to risk across states of
\cref{ch:b2:expected-utility}, measured by $\gamma$, and the reluctance to let consumption change
across time, measured by $1/\gamma$. A household that dislikes gambles must, in this model, also
dislike uneven consumption over time. Nothing in the evidence forces the two to be linked; the model
links them for tractability.

\begin{omfigure}
\begin{tikzpicture}[econ]
\begin{axis}[omrecord, grid=none, xmin=0, xmax=90, ymin=0, ymax=90, xtick={20,39.56,77.14}, xticklabels={20,
    39.6, 77.1}, ytick={39.46,60,81}, yticklabels={39.5,60,81},
    xlabel={Consumption now, $c_0$ (thousands)}, ylabel={Consumption next year, $c_1$ (thousands)},
    width=0.62\linewidth, height=0.55\linewidth]
  \addplot[omDef, very thick, domain=0:77.14, samples=2] {(77.142857 - x)*1.05};
  \addplot[black, thick, domain=24:88, samples=80] {39.461538*(39.560440/x)^(1/0.95)};
  \draw[guide] (axis cs:20,0) -- (axis cs:20,60) -- (axis cs:0,60);
  \draw[guide] (axis cs:39.56,0) -- (axis cs:39.56,39.46) -- (axis cs:0,39.46);
  \node[eqpoint] at (axis cs:20,60) {};
  \node[anchor=south west, font=\small] at (axis cs:20,60) {$E$};
  \node[eqpoint] at (axis cs:39.56,39.46) {};
  \node[anchor=south west, font=\small] at (axis cs:39.56,39.46) {$A$};
  \draw[->, omDef] (axis cs:66,52) -- (axis cs:66.5,12);
  \node[anchor=south, font=\footnotesize, omDef] at (axis cs:66,52) {budget line, slope $-(1+r)$};
\end{axis}
\end{tikzpicture}

{\small Two periods, $u = \ln$, $\beta = 0.95$, $r = 5$ per cent. A student earns 20 now and 60 next
year (thousands of dalens; point $E$). Lifetime resources are $20 + 60/1.05 = 77.1$. The best plan
$A$ is where an indifference curve touches the budget line: there the slope of the curve,
$-u'(c_0)/(\beta u'(c_1))$, equals $-(1 + r)$, which is the \omterm{def:b2:intertemporal-choice:euler-equation}{Euler equation}. Since $\beta(1 + r) =
0.9975$, the plan is almost flat; the student borrows 19.6 now.}
\end{omfigure}

\section{Consumption smoothing}\label{sec:b2:intertemporal-choice:smoothing}

\begin{definition}[Consumption smoothing]\label{def:b2:intertemporal-choice:smoothing}
A household practises \emph{consumption smoothing}\index{consumption smoothing} when it uses borrowing
and saving to make its consumption path less variable than its income path: it consumes more than its
income when income is low and less when it is high.
\end{definition}

\begin{proposition}[Flat consumption]\label{prop:b2:intertemporal-choice:flat}
If $\beta(1 + r) = 1$, the optimal plan is flat, whatever the income profile and whatever the
(strictly concave) utility:
\[
c_t = \bar c = \frac{\mathcal{M}}{\sum_{t=0}^{T-1}(1 + r)^{-t}} \quad\text{for every } t.
\]
A change in income in a single period changes consumption in every period, by the same amount, and by
much less than the change in income.
\end{proposition}

\begin{proof}
With $\beta(1 + r) = 1$ the \omterm{def:b2:intertemporal-choice:euler-equation}{Euler equation} is $u'(c_t) = u'(c_{t+1})$; $u'$ is strictly decreasing, so
$c_t = c_{t+1}$. The budget constraint then gives $\bar c\sum_t (1 + r)^{-t} = \mathcal{M}$. A rise
$\Delta$ in $m_s$ raises $\mathcal{M}$ by $\Delta/(1 + r)^{s}$ and so every $c_t$ by that amount
divided by $\sum_t (1 + r)^{-t}$, which is larger than one when $T > 1$.
\end{proof}

The listing below works out a whole life. Each period is a decade, so the interest rate over a period
is 48 per cent (about 4 per cent a year); income rises to a peak in the fourth decade and falls in the
sixth. The program finds the consumption path from the \omterm{def:b2:intertemporal-choice:euler-equation}{Euler equation} and the budget, then follows
assets decade by decade.

\omcode{code/bachelor-2/03-intertemporal-choice/listings/life_plan.py}{1}{26}{A life of six decades:
the growth of consumption from the \omterm{def:b2:intertemporal-choice:euler-equation}{Euler equation}, its level from the budget, then the assets.}
\label{lst:b2:intertemporal-choice:life-plan}

\omoutput{code/bachelor-2/03-intertemporal-choice/outputs/life_plan.txt}{A patient household and an
impatient one, on the same income.}\label{out:b2:intertemporal-choice:life-plan}

With $\beta(1 + r) = 1$, consumption is flat at 28.3 thousand dalens a year. The household borrows in
its first two decades, repays its debt in the third, saves in the fourth and fifth, and spends its
savings in the last, ending with nothing. (The model treats a decade as one period, so assets are
printed per year of the decade, like income.) The impatient household, with $\beta(1 + r) = 0.74$, borrows more when young and lets
its consumption fall by 14 per cent a decade, $0.74^{1/2} = 0.86$. Neither path follows the hump of
income. Both are what the model predicts for a household that can borrow at the market rate; the next
two sections say what changes when it cannot, or when the interest rate moves.

\begin{measure}[Spending by age in an expenditure survey]
An expenditure survey interviews a sample of households (``consumer units'') about their spending and
income, and the agency publishes averages by the age of the \emph{reference person}, the member who
owns or rents the home. Reading the opening chart as a life cycle needs three cautions. It is a
cross-section: in 2023 the households aged 75 were born some thirty years before those aged 45 and
earned their incomes in a different economy, so the chart mixes age with the date of birth. Households change size: an average unit
% ledger: b2:bls-ce-age-35-44, b2:bls-ce-age-75p
has 3.4 people at 35--44 and 1.6 at 75 and older, so spending per person is much flatter than spending
per household: from the youngest group to its peak (55--64) it rises 1.48 times, against 1.85 for
% figdata: ce_summary (peak_pp_spending_ratio 1.48)
the unit to its own peak (45--54). And the survey's ``expenditures'' include pension and Social Security
contributions, which the chart removes: pension contributions are saving and Social Security contributions a
tax, and neither is consumption.
\end{measure}

\begin{model}{Exponential discounting and consumption smoothing}
\assumes A household with time-separable utility and a constant \omterm{def:b2:intertemporal-choice:discount-factor}{discount factor} plans its whole life,
knows its income profile, and borrows and lends at one interest rate.
\predicts Consumption grows at $(\beta(1 + r))^{1/\gamma}$ whatever the income profile; it is flat when
$\beta(1 + r) = 1$; predictable changes in income do not move consumption when they happen; plans,
once made, are carried out.
\tested Expenditure surveys by age (this chapter's opening chart): spending varies less than income,
and the oldest households spend more than their income.
\breaks Spending still rises and falls with income over the life cycle, more than smoothing allows,
part of which household size and cohorts explain; households that cannot borrow
(\cref{sec:b2:intertemporal-choice:borrowing}); preference reversals and the demand for commitment
(\cref{sec:b2:intertemporal-choice:consistency}).
\end{model}

\section{The interest rate}\label{sec:b2:intertemporal-choice:rate}

The Year 1 volume showed with two periods that a rise in the interest rate has a substitution effect
(consumption today becomes dearer, so save more) and an income effect (a saver is richer, so consume
more today). With CRRA utility the balance has a clean answer.

\begin{proposition}[Saving and the interest rate]\label{prop:b2:intertemporal-choice:rate}
Two periods, all income $m_0$ in period 0, CRRA utility. First-period consumption is
\[
c_0 = \frac{m_0}{1 + \beta^{1/\gamma}(1 + r)^{1/\gamma - 1}} .
\]
A rise in $r$ lowers $c_0$ (raises saving) if $\gamma < 1$, leaves it unchanged if $\gamma = 1$, and
raises $c_0$ (lowers saving) if $\gamma > 1$.
\end{proposition}

\begin{proof}
The budget is $c_1 = (1 + r)(m_0 - c_0)$ and the \omterm{def:b2:intertemporal-choice:euler-equation}{Euler equation} $c_1 = (\beta(1 + r))^{1/\gamma}c_0$.
Equating, $m_0 - c_0 = c_0\,\beta^{1/\gamma}(1 + r)^{1/\gamma - 1}$, which gives $c_0$. The term $(1 +
r)^{1/\gamma - 1}$ rises with $r$ when $1/\gamma > 1$, is constant when $\gamma = 1$, and falls when
$1/\gamma < 1$; $c_0$ moves the other way.
\end{proof}

The \omterm{def:b2:intertemporal-choice:eis}{elasticity of intertemporal substitution} measures the substitution effect: when it exceeds one,
it wins. With logarithmic utility the two effects cancel exactly, and the household consumes the
share $1/(1 + \beta)$ of its wealth whatever the rate. When income also arrives later, a third effect
appears: a higher rate lowers the present value of future income, the household's \emph{human wealth},
which makes it poorer and pushes saving up. So a rise in the interest rate raises the saving of a young
household with most of its income ahead of it more surely than the saving of a retiree living on its
wealth. How large the EIS is, and so how much saving responds to the interest rate, is an empirical
question that the aggregate data leave open; \cref{ch:b2:consumption} returns to it.

\section{Borrowing limits}\label{sec:b2:intertemporal-choice:borrowing}

Smoothing a rising income profile needs borrowing against future income, which lenders may refuse:
future earnings cannot be pledged as collateral, and a borrower may not repay.

\begin{definition}[Borrowing constraint]\label{def:b2:intertemporal-choice:borrowing-constraint}
A \emph{borrowing constraint}\index{borrowing constraint} is a limit $b_t \ge -\bar B$ on the debt a
household may hold at the end of each period, with $\bar B \ge 0$; $\bar B = 0$ forbids borrowing.
\end{definition}

\begin{proposition}[Euler equation with a borrowing limit]\label{prop:b2:intertemporal-choice:constrained}
Under a \omterm{def:b2:intertemporal-choice:borrowing-constraint}{borrowing constraint} the optimal plan satisfies, in each period,
\[
u'(c_t) \ge \beta(1 + r)\, u'(c_{t+1}),
\]
with equality whenever the limit does not bind at the end of period $t$ ($b_t > -\bar B$). When it
binds, consumption is lower today than the \omterm{def:b2:intertemporal-choice:euler-equation}{Euler equation} wants: $c_t = m_t + (1 + r)b_{t-1} + \bar B$,
and consumption moves with current income.
\end{proposition}

\begin{proof}
Add to the Lagrangian of \cref{prop:b2:intertemporal-choice:euler} the multipliers $\mu_t \ge 0$ of the
period constraints $b_t + \bar B \ge 0$, written in terms of $c_0, \dots, c_t$. Raising $c_t$ by one and
lowering $c_{t+1}$ by $1 + r$ leaves every later asset unchanged and lowers only $b_t$, by one. If
$b_t > -\bar B$, this swap and its reverse are both feasible, so neither may raise utility: equality. If
$b_t = -\bar B$, only the reverse swap (less now, more later) is feasible, and it may not raise
utility: $\beta^{t}u'(c_t) \ge \beta^{t+1}(1 + r)u'(c_{t+1})$.
\end{proof}

\begin{omfigure}
\begin{tikzpicture}[econ]
\begin{axis}[omrecord, grid=none, xmin=0, xmax=90, ymin=0, ymax=90, xtick={20,39.56}, xticklabels={20,
    39.6}, ytick={60,81}, xlabel={Consumption now, $c_0$ (thousands)},
    ylabel={Consumption next year, $c_1$ (thousands)}, width=0.62\linewidth, height=0.55\linewidth]
  \addplot[omDef, very thick, domain=0:20, samples=2] {(77.142857 - x)*1.05};
  \draw[omDef, very thick] (axis cs:20,60) -- (axis cs:20,0);
  \addplot[omDef, thick, dashed, domain=20:77.14, samples=2] {(77.142857 - x)*1.05};
  \addplot[black, thick, domain=14.5:60, samples=80] {60*(20/x)^(1/0.95)};
  \addplot[black!50, thin, domain=24:88, samples=80] {39.461538*(39.560440/x)^(1/0.95)};
  \node[eqpoint] at (axis cs:20,60) {};
  \node[anchor=south west, font=\small] at (axis cs:20,60) {$E$};
  \node[circle, draw=black!60, inner sep=0pt, minimum size=4pt] at (axis cs:39.56,39.46) {};
  \node[anchor=south west, font=\small, black!60] at (axis cs:39.56,39.46) {$A$};
  \node[anchor=west, font=\footnotesize, omDef] at (axis cs:21,8) {no borrowing: $c_0 \le m_0$};
\end{axis}
\end{tikzpicture}

{\small The same student when nobody lends to it ($\bar B = 0$). The feasible set loses the dashed
part of the budget line, and the best plan is the corner $E$: the student consumes its income. The
indifference curve through $E$ is steeper than the budget line there ($u'(c_0)/(\beta u'(c_1)) = 60/(0.95
\times 20) = 3.2$, against $1.05$): the student would pay far more than the market rate to borrow. $A$
is the plan it cannot reach.}
\end{omfigure}

A binding limit changes the predictions. Consumption tracks current income for households at the
limit, mostly young ones whose income will rise; a predictable rise in income raises their consumption
when it arrives, not before; and a loan, a transfer or a tax refund that arrives early raises their
spending one for one. The limit also has a cost, which \cref{exo:b2:intertemporal-choice:12} measures:
the share of lifetime resources the household would give up to be rid of it.

Brenland's household survey shows the extreme case. Grouping its 2,000 households by the age of the
head, consumption is between 92 and 97 per cent of disposable income in every age band: it follows the
% figdata: brenland_age (ratio 0.924 to 0.973)
hump of income almost exactly. Brenland's invented households behave as if every one of them were at
its borrowing limit, or did not plan at all.

\begin{omchart}{Brenland's household survey, 2025 (invented): mean disposable income and mean
  consumption by age band of the head, thousands of dalens a year; 287 to 370 households per band.}{Brenland
  household survey (invented data, \texttt{households.csv})}
\begin{tikzpicture}
\begin{axis}[omrecord, xmin=20, xmax=80, ymin=0, ymax=35, xtick={25,34.5,44.5,54.5,64.5,74.5},
    xticklabels={21--29,30--39,40--49,50--59,60--69,70--79}, xlabel={Age of the head},
    ylabel={Thousands of dalens a year}, legend style={at={(0.5,0.04)}, anchor=south}, legend columns=2,
    height=0.45\linewidth, width=0.8\linewidth]
  \addplot[omDef, very thick, mark=*, mark size=1.8pt] table[x=mid, y=disposable]
    {figdata/bachelor-2/03-intertemporal-choice/out/brenland_age.csv};
  \addplot[omThm, very thick, mark=square*, mark size=1.8pt] table[x=mid, y=consumption]
    {figdata/bachelor-2/03-intertemporal-choice/out/brenland_age.csv};
  \legend{Disposable income, Consumption}
\end{axis}
\end{tikzpicture}
\end{omchart}

\section{Time consistency}\label{sec:b2:intertemporal-choice:consistency}

A plan is made today but carried out later, by the same person at later dates. Will that later person
still want it?

\begin{definition}[Time consistency]\label{def:b2:intertemporal-choice:time-consistency}
Preferences are \emph{time consistent}\index{time consistency} if a plan that is optimal at date 0 has
the property that, at every later date $t$, its remainder is optimal for the household at $t$, given
what it has done up to $t$. A household with time-consistent preferences never wants to revise its
plan when nothing new has been learnt.
\end{definition}

\begin{proposition}[Exponential discounting is time consistent]\label{prop:b2:intertemporal-choice:consistent}
With exponential discounting and strictly concave $u$, the remainder from date $t$ of the date-0
optimal plan is the date-$t$ optimal plan, given assets $b_{t-1}$.
\end{proposition}

\begin{proof}
Write $U_0 = \sum_{k<t}\beta^{k}u(c_k) + \beta^{t}U_t$, with $U_t = \sum_{k\ge t}\beta^{k-t}u(c_k)$
the date-$t$ objective. Let $(c^*_k)$ be the date-0 optimum and $b^*_{t-1}$ the assets it leaves at
$t$. The plans from $t$ on that are feasible with $b^*_{t-1}$ are the same for both selves. If some
feasible remainder gave a higher $U_t$ than $(c^*_k)_{k\ge t}$, then splicing it onto
$(c^*_k)_{k<t}$ would give a higher $U_0$, since $\beta^{t} > 0$, contradicting optimality at date 0.
So the remainder maximises $U_t$; the maximum is unique because $U_t$ is strictly concave.
\end{proof}

The proof works because the weight the date-0 self puts on period $t + 1$ relative to period $t$,
$\beta^{t+1}/\beta^{t} = \beta$, is the same as the weight the date-$t$ self puts on them. Only
exponential discounting has that property (\cref{exo:b2:intertemporal-choice:10}). The simplest
departure keeps $\beta$ between any two future periods but adds an extra discount between now and
everything later.

\begin{definition}[Quasi-hyperbolic discounting]\label{def:b2:intertemporal-choice:quasi-hyperbolic}
A household with \emph{quasi-hyperbolic discounting}\index{quasi-hyperbolic discounting} ranks plans at
date $t$ by
\[
U_t = u(c_t) + \varphi\sum_{k=1}^{T-1-t}\beta^{k}u(c_{t+k}), \qquad 0 < \varphi \le 1,\ 0 < \beta \le 1.
\]
With $\varphi = 1$ this is exponential discounting.
\end{definition}

\begin{remark}
Much of the literature writes this model with $\beta$ for the present-bias factor and $\delta$ for the
long-run \omterm{def:b2:intertemporal-choice:discount-factor}{discount factor}, and calls it the $\beta$--$\delta$ model. This book keeps $\beta$ for the
\omterm{def:b2:intertemporal-choice:discount-factor}{discount factor} and $\delta$ for depreciation, so the present-bias factor is $\varphi$.
\end{remark}

\begin{definition}[Present bias]\label{def:b2:intertemporal-choice:present-bias}
A household has \emph{present bias}\index{present bias} if it weights the present more heavily, relative
to the future, than its own discounting between any two future dates implies; in the quasi-hyperbolic
model, if $\varphi < 1$.
\end{definition}

\begin{proposition}[Present bias is time inconsistent]\label{prop:b2:intertemporal-choice:inconsistent}
With $\varphi < 1$ and an interior optimum, the date-0 plan for periods 1 and 2 satisfies $u'(c_1) =
\beta(1 + r)u'(c_2)$, but the household at date 1 chooses $u'(c_1) = \varphi\beta(1 + r)u'(c_2)$: it
consumes more at date 1 than it planned.
\end{proposition}

\begin{proof}
At date 0 the weights on periods 1 and 2 are $\varphi\beta$ and $\varphi\beta^2$, whose ratio is
$\beta$; the Euler argument of \cref{prop:b2:intertemporal-choice:euler} between those two periods
gives the first equation. At date 1 the weights are $1$ and $\varphi\beta$, whose ratio is $\varphi\beta
< \beta$. Since $u'$ is decreasing, the second equation, with a smaller right-hand side for a given
$c_2$, needs a larger $c_1$ out of the same resources.
\end{proof}

The clearest symptom is a \emph{preference reversal}. Offer \dal{100} on day 10 or \dal{110} on day
11, and take $\beta = 1$ between days for simplicity, with $\varphi = 0.7$. Seen from day 0, both dates
are in the future and carry the same weight $0.7$: the later, larger amount wins, 77 against 70. On
day 10, the first amount is now and the second is not: 100 against 77, and the sooner amount wins. The
household chooses, then changes its mind with no new information.

\begin{omfigure}
\begin{tikzpicture}[econ, x=1cm, y=0.028cm]
  \foreach \xo/\title/\a/\b in {0/{Seen from day 0}/70/77, 6.4/{Seen on day 10}/100/77} {
    \begin{scope}[xshift=\xo cm]
      \draw[econ axes, <-] (0,115) -- (0,0) -- (5,0);
      \fill[omDef!60] (1,0) rectangle (2,\a);
      \fill[omThm!60] (3,0) rectangle (4,\b);
      \node[above, font=\footnotesize] at (1.5,\a) {\a};
      \node[above, font=\footnotesize] at (3.5,\b) {\b};
      \node[below, font=\footnotesize, align=center] at (1.5,0) {\dal{100}\\day 10};
      \node[below, font=\footnotesize, align=center] at (3.5,0) {\dal{110}\\day 11};
      \node[font=\small] at (2.5,128) {\title};
    \end{scope}
  }
  \node[font=\footnotesize, anchor=west] at (0.1,108) {value now};
\end{tikzpicture}

{\small A preference reversal with $\varphi = 0.7$ and $\beta = 1$ between days. From day 0 the two
offers are both discounted by $0.7$, and \dal{110} on day 11 is worth more (77 against 70). On day 10
the first offer is no longer discounted, and \dal{100} now beats \dal{110} tomorrow (100 against 77).
Under exponential discounting the comparison is the same on both days.}
\end{omfigure}

What the household does about its future selves depends on whether it knows they will deviate. A
\emph{naive} household believes its future selves will follow today's plan; a \emph{sophisticated}
one predicts their choices correctly and plans around them. Only the sophisticated household has a
reason to tie its own hands.

\begin{definition}[Commitment device]\label{def:b2:intertemporal-choice:commitment-device}
A \emph{commitment device}\index{commitment device} is an arrangement, chosen freely, that removes some
of a person's future options or makes them costlier: a savings account that cannot be touched before a
date, a contract that charges for missed sessions, a deadline set in advance.
\end{definition}

\begin{proposition}[Who pays for commitment]\label{prop:b2:intertemporal-choice:commitment}
A household with exponential discounting never pays a positive price for a \omterm{def:b2:intertemporal-choice:commitment-device}{commitment device} that
removes options. A sophisticated household with \omterm{def:b2:intertemporal-choice:present-bias}{present bias} may.
\end{proposition}

\begin{proof}
By \cref{prop:b2:intertemporal-choice:consistent} the exponential household's later selves carry out
its date-0 plan, which is feasible with or without the device; removing options can only shrink what
is feasible, and a price lowers resources, so its utility falls. For the second claim it is enough to
give a case, which \cref{exo:b2:intertemporal-choice:11} and the weekend problem work out: the
date-1 self consumes more than the date-0 self wants (\cref{prop:b2:intertemporal-choice:inconsistent}),
and a device that forces the date-0 plan raises the date-0 self's utility.
\end{proof}

So a demand for commitment at a positive price refutes exponential discounting, whatever the \omterm{def:b2:intertemporal-choice:discount-factor}{discount
factor}; the next section reports two such observations.

\section{Measured present bias}\label{sec:b2:intertemporal-choice:measured}

\begin{record}[Paying not to go to the gym]
DellaVigna and Malmendier (2006) followed the contract choices and daily attendance of 7,752 members of
% ledger: b2:dm2006
three American health clubs over three years. Members who chose a flat monthly fee of over \usd{70}
attended 4.3 times a month on average: more than \usd{17} per visit, when a 10-visit pass cost
\usd{10} per visit. On average they forwent \usd{600} over their membership. Monthly members were also
17 per cent more likely than annual members to stay enrolled beyond a year, although the monthly
contract costs more precisely for the option to cancel. The authors' leading explanation is
overconfidence about future self-control: members who expect to go often, and to cancel promptly,
and do neither. In the terms of this chapter, a naive present-biased member does what the data show.
\end{record}

\begin{record}[Effort today or tomorrow]
Augenblick, Niederle and Sprenger (2013) had 102 participants allocate a number of dull tasks between two
% ledger: b2:ans2013-design
dates, over seven weeks: once a week in advance, and again on the day the first date arrived. On the
day, they allocated about 9 per cent fewer tasks to the present than they had planned a week before.
% ledger: b2:ans2013-beta
The estimated present-bias factor is $0.900$ (standard error $0.037$) for effort, while the same
participants choosing between payments of money show $0.974$ (standard error $0.009$), close to no \omterm{def:b2:intertemporal-choice:present-bias}{present
bias}. When offered, for free, a device that made their advance allocation more likely to count, 59 per
% ledger: b2:ans2013-commit
cent of them took it, and those whose allocations had shown \omterm{def:b2:intertemporal-choice:present-bias}{present bias} were the ones who chose it.
But almost nobody would pay for it: at a price of 25 cents, 91 per cent preferred flexibility.
\end{record}

\begin{model}{Quasi-hyperbolic discounting}
\assumes A household weights every future period by an extra factor $\varphi < 1$ relative to the
present, and discounts exponentially between future periods; it may be naive or sophisticated about
its future selves.
\predicts Preference reversals as a reward approaches; plans to save, exercise or work that are not
carried out; a demand for \omterm{def:b2:intertemporal-choice:commitment-device}{commitment devices} among the sophisticated, even at a price.
\tested Gym contracts and attendance (DellaVigna and Malmendier 2006); effort allocated in advance
and on the day, and the take-up of commitment (Augenblick, Niederle and Sprenger 2013).
\breaks The estimated $\varphi$ depends on what is being chosen: well below one for effort, close to
one for money in the same people; and few people pay a positive price for commitment in the
effort experiment, so the size of the demand is small.
\end{model}

The two models are less opposed than they look. Exponential discounting describes the long-run shape
of plans, and its predictions about consumption growth and the interest rate are what the rest of the
book uses, in the asset-pricing chapter and in the macroeconomic models. \omterm{def:b2:intertemporal-choice:present-bias}{Present bias} describes the
gap between plans and their execution, which matters most where a cost comes now and a benefit later.
Neither model says which choices people \emph{should} make; each predicts which ones they \emph{will}
make, and the records above say where each prediction holds.

\section{Exercises}\label{sec:b2:intertemporal-choice:exercises}

\begin{exercise}[$\star$]\label{exo:b2:intertemporal-choice:1}
An apprentice will earn \dal{30000} this year, \dal{40000} next year and \dal{50000} the year after.
The interest rate is 4 per cent. Find the present value of these earnings.
\end{exercise}

\begin{exercise}[$\star$]\label{exo:b2:intertemporal-choice:2}
A household lives two periods, earns $m_0 = \num{40000}$ and $m_1 = \num{42000}$ dalens, has $u = \ln$,
$\beta = 0.95$ and borrows or lends at $r = 5$ per cent. Find $c_0$, $c_1$ and its saving in period 0.
\end{exercise}

\begin{exercise}[$\star$]\label{exo:b2:intertemporal-choice:3}
Read the opening chart's data (the ledger rows, or the table of the chart's source). (a) What share of
its income after taxes does the 25--34 group spend, and the group aged 75 and older? (b) Spending per
person divides by the average number of people in the unit: find it for the youngest group and for the
55--64 group. (c) Which comparison, per unit or per person, is closer to what smoothing predicts?
\end{exercise}

\begin{exercise}[$\star$]\label{exo:b2:intertemporal-choice:4}
Name two observations that would refute exponential discounting, whatever the \omterm{def:b2:intertemporal-choice:discount-factor}{discount factor}. Then
say why a household that consumes more when young than when old does not refute it.
\end{exercise}

\begin{exercise}[$\star\star$]\label{exo:b2:intertemporal-choice:5}
A household has CRRA utility with $\gamma = 2$, a \omterm{def:b2:intertemporal-choice:time-preference}{rate of time preference} $\rho = 2$ per cent, and
faces $r = 4$ per cent. Find the growth rate of its consumption, exactly and with the approximation of
\cref{cor:b2:intertemporal-choice:crra-growth}. Do the same with $\gamma = 0.5$. Which household has
the larger EIS?
\end{exercise}

\begin{exercise}[$\star\star$]\label{exo:b2:intertemporal-choice:6}
A household has \dal{100000} now and nothing next period, $\beta = 0.96$. The interest rate rises from
3 to 6 per cent. Predict the sign of the change in its saving for $\gamma = 2$, $1$ and $0.5$, then
compute $c_0$ before and after in each case. Which effect wins in each?
\end{exercise}

\begin{exercise}[$\star\star$]\label{exo:b2:intertemporal-choice:7}
The student of the figures earns \dal{20000} now and \dal{60000} next year, with $u = \ln$, $\beta =
0.95$, $r = 5$ per cent. (a) Find its unconstrained plan. (b) Nobody lends to it. Show that the limit
binds by comparing $u'(c_0)/u'(c_1)$ at the endowment with $\beta(1 + r)$. (c) What does it consume?
\end{exercise}

\begin{exercise}[$\star\star$]\label{exo:b2:intertemporal-choice:8}
Data and Python. Load Brenland's survey with \texttt{from country import load};
\texttt{load("households")}. Group the households by the age of the head in bands of ten years from 21,
and compute mean disposable income, mean consumption and their ratio in each band. Fit consumption on
disposable income across households by least squares. Does consumption track income over the life
cycle? Compare the highest band mean with the lowest, for income and for consumption, and say what
smoothing would predict.
\end{exercise}

\begin{exercise}[$\star\star$]\label{exo:b2:intertemporal-choice:9}
A bank offers a savings account that cannot be touched for a year, for a fee of \dal{20}. (a) Show that
a household with exponential discounting never opens it, whatever $\beta$. (b) A naive present-biased
household does not open it either. Why? (c) Which of the two records of
\cref{sec:b2:intertemporal-choice:measured} shows a demand for commitment, and at what price?
\end{exercise}

\begin{exercise}[$\star\star\star$]\label{exo:b2:intertemporal-choice:10}
Let a household rank plans at date $t$ by $\sum_{k\ge 0} D(k)\,u(c_{t+k})$, with $D(0) = 1$ and $D(k) >
0$. Show that its preferences are \omterm{def:b2:intertemporal-choice:time-consistency}{time consistent} for every utility $u$ and every budget only if $D(k)
= D(1)^{k}$: exponential discounting is the only time-consistent discounting of this form. (Compare
the marginal rate of substitution between dates $t + k$ and $t + k + 1$ seen from date 0 and from date
$t$.)
\end{exercise}

\begin{exercise}[$\star\star\star$]\label{exo:b2:intertemporal-choice:11}
A household has \dal{90000} to eat over three periods, no interest ($r = 0$), $u = \ln$, $\beta = 1$
and $\varphi = 0.5$. (a) Find what the date-1 self consumes out of what is left, $W_1$. (b) Find the
sophisticated date-0 self's $c_0$, and the whole path. (c) What split of $W_1$ would the date-0 self
like? (d) Show that a naive date-0 self chooses the same $c_0$ with logarithmic utility.
\end{exercise}

\begin{exercise}[$\star\star\star$]\label{exo:b2:intertemporal-choice:12}
For the student of \cref{exo:b2:intertemporal-choice:7}, find the share $\lambda$ of its lifetime
resources that would give it, without a borrowing limit, the same utility as its constrained plan. (With
two periods and $u = \ln$, lifetime utility is $(1 + \beta)\ln(\text{resources}) + $ a constant.) Express
the cost of the limit in dalens of lifetime resources.
\end{exercise}

\section{Problem: An engineer in Brenmoor plans her life}\label{sec:b2:intertemporal-choice:problem}

\begin{problem}[{Weekend problem --- an engineer in Brenmoor, the capital of Brenland, plans her
consumption over three stretches of her life: smoothing, a borrowing limit, and a present bias that a
commitment savings account can correct}]
\label{pb:b2:intertemporal-choice:1}
An engineer lives three periods: young, middle-aged and retired. She earns $m_0 = \num{24000}$, $m_1 =
\num{60000}$ and $m_2 = \num{6000}$ dalens per period, has $u = \ln$ and $\beta = 0.8$, and borrows or
lends at $r = 25$ per cent per period. Round amounts to the nearest dalen and utilities to four
decimals.

\medskip
\noindent\textbf{Part I --- Budget and present value.}
\begin{enumerate}
  \item Write her budget period by period, with assets $b_0$ and $b_1$ and $b_2 = 0$, and derive the
        intertemporal budget constraint.
  \item Find the present value $\mathcal{M}$ of her income.
  \item Find her \omterm{def:b2:intertemporal-choice:time-preference}{rate of time preference}, and compare it with $r$.
  \item If she consumed her income each period, what would her lifetime utility be? Find the constant
        consumption $\bar c_h$ that gives the same utility.
\end{enumerate}

\medskip
\noindent\textbf{Part II --- Euler equation and smoothing.}
\begin{enumerate}[resume]
  \item Write the two \omterm{def:b2:intertemporal-choice:euler-equation}{Euler equations} and show that her plan is flat.
  \item Find her consumption.
  \item Find $b_0$ and $b_1$, and check that $b_2 = 0$.
  \item By how much per period does smoothing beat consuming her income, in constant consumption?
  \item The interest rate rises to 50 per cent. Find the growth factor of her consumption, the new
        $\mathcal{M}$, and $c_0$, $c_1$, $c_2$.
  \item Does she borrow more or less when young after the rise? Show that with $u = \ln$, $c_0 =
        \mathcal{M}/(1 + \beta + \beta^2)$ at every rate, and say which effect of the rate remains.
  \item With $r = 25$ per cent again, would her plan change if $\gamma$ were 2 instead of 1? Why?
\end{enumerate}

\medskip
\noindent\textbf{Part III --- A borrowing limit.}
\begin{enumerate}[resume]
  \item Banks do not lend to the young: $b_0 \ge 0$. Does the limit bind?
  \item Find her constrained plan.
  \item Check the Euler inequality of \cref{prop:b2:intertemporal-choice:constrained} between periods
        0 and 1, and the \omterm{def:b2:intertemporal-choice:euler-equation}{Euler equation} between 1 and 2.
  \item Compute her lifetime utility with and without the limit.
  \item Find the constant consumption $\bar c$ that gives her the constrained utility, the cost of the
        limit per period, and its cost in dalens of lifetime resources.
  \item A lender offers to lift the limit for a fee of \dal{1000} in present value. Does she accept?
\end{enumerate}

\medskip
\noindent\textbf{Part IV --- Present bias and a commitment savings account.}
\begin{enumerate}[resume]
  \item She now has $\varphi = 0.7$, and no borrowing limit. Write the weights her date-0 and date-1
        selves put on the three periods.
  \item What split of the resources $W_1$ left at date 1 would her date-0 self like?
  \item What split does her date-1 self choose?
  \item She is sophisticated. Find $c_0$, $W_1$, $c_1$ and $c_2$. (Use $u = \ln$: the date-0 self
        values $W_1$ through $\varphi\beta(1 + \beta)\ln W_1$ plus a constant.)
  \item Would a naive engineer choose a different $c_0$? Explain.
  \item A commitment savings account lets her fix $c_1$ and $c_2$ at date 0. What does she choose, and
        how many dalens of $W_1$ would her date-0 self give up for the account?
  \item Compare the three plans --- exponential ($\varphi = 1$), present-biased without the account,
        present-biased with it --- by the exponential utility $\ln c_0 + \beta\ln c_1 + \beta^2\ln c_2$.
        What does the account fix, and what does it leave?
  \item State the result: the cost of the borrowing limit and the value of the commitment account, in
        dalens.
\end{enumerate}
\end{problem}
